%% Determines the type of document (including standard settings for layout).
\documentclass[letterpaper]{article}



%% Package to control the size of a page: the standard margins used by LaTeX are
%% very wide!
\usepackage[margin=1in]{geometry}


%% Packages add functionality to the document. The AMS packages are standard
%% packages to support various mathematical notations. AMS stands for American
%% Mathematical Society, the organization that maintains these packages.
\usepackage{amsmath,amsthm,amssymb}

%% Theorems-like environments using functionality provided by amsthm.
\theoremstyle{plain}
\newtheorem{theorem}{Theorem}[section]
    %% [section] at the end specifies that theorems should be numbered
    %% per-section: Section x starts with theorem-like x.1 and so on, ...
\newtheorem{proposition}[theorem]{Proposition}
    %% [theorem] in the middle specifies: use the same counter as the theorem
    %% environment: here we number all theorem-like environments consecutively.
\newtheorem{corollary}[theorem]{Corollary}
\newtheorem{lemma}[theorem]{Lemma}
\theoremstyle{definition}
\newtheorem{definition}[theorem]{Definition}
\theoremstyle{remark}
\newtheorem{example}[theorem]{Example}
\newtheorem{remark}[theorem]{Remark}


%% Support for nicely formatted tables.
\usepackage{booktabs}


%% Support for colors & colors in tables.
\usepackage[table]{xcolor}

% Seven colors safe for use color blindness.
% Colors taken from doi:10.1038/nmeth.1618.
\definecolor{cbOrange}{RGB}{230,159,0}
\definecolor{cbSkyBlue}{RGB}{86,180,233}
\definecolor{cbBluishGreen}{RGB}{0,158,115}
\definecolor{cbYellow}{RGB}{240,228,66}
\definecolor{cbBlue}{RGB}{0,114,178}
\definecolor{cbVermillion}{RGB}{213,94,0}
\definecolor{cbReddischPurple}{RGB}{204,121,167}


%% Notation used in this document.
\newcommand{\n}{\mathbf{n}} %% Num. Replicas.
\newcommand{\f}{\mathbf{f}} %% Num. Faulty Replicas.

%% Misc. Math notation.
\newcommand{\BigO}{\mathcal{O}}
\newcommand{\abs}[1]{\lvert#1\rvert}
\newcommand{\AName}[1]{\textsc{#1}}
\newcommand{\Var}[1]{\texttt{#1}}


%% Algorithms.
\usepackage{algorithm}
\usepackage[noend]{algorithmic}
\newcommand{\GETS}{:=}


%% Formatting SI-units.
\usepackage{siunitx}
\sisetup{per-mode=symbol}


%% TikZ: for creating figures.
\usepackage{tikz}

%% Configuration for figures: Nicer arrows.
\usetikzlibrary{arrows.meta}
\tikzset{>=Stealth}


%% pgfplots: drawing plots using TikZ.
\usepackage{pgfplots}
%% Configuration for plots: Use color-blind friendly colors.
\pgfplotscreateplotcyclelist{cbSafeList}{
    very thick,solid,cbOrange,every mark/.append style={solid},mark=*\\
    very thick,solid,cbSkyBlue,every mark/.append style={solid},mark=*\\
    very thick,solid,cbBluishGreen,every mark/.append style={solid},mark=*\\
    very thick,solid,cbYellow,every mark/.append style={solid},mark=*\\
    very thick,solid,cbBlue,every mark/.append style={solid},mark=*\\
    very thick,solid,cbVermillion,every mark/.append style={solid},mark=*\\
    very thick,solid,cbReddischPurple,every mark/.append style={solid},mark=*\\
    very thick,solid,black,every mark/.append style={solid},mark=*\\
}
\pgfplotsset{
    legend style={font=\small},
    compat=1.16,
    width=260pt,
    height=140pt,
    legend cell align=left,
    xlabel near ticks,
    ylabel near ticks,
    every axis/.append style={
        cycle list name=cbSafeList,
        ymin=0,
        enlargelimits=0.05,
        mark size=1pt,
        ylabel style={align=center},
        xlabel style={align=center},
        title style={align=center}
    }
}

%% PgfplotsTable: loading data files to use with pgfplots.
\usepackage{pgfplotstable}


%% Support for hyperlinks and urls. The setting ``colorlinks'' sets how links
%% are shown in the document (with a color, without underline). We put hyperref
%% last---it has a tendency to break other packages when loaded before them.
\usepackage[colorlinks]{hyperref}
\usepackage{graphicx}
\usepackage{pgffor}
\usepackage{caption}
\usepackage{tabularx}
\usepackage{enumitem}
\usepackage{fancyhdr}
\usepackage{comment}
\usepackage{titlesec}


\begin{document}

%% First the meta-data of the doucment: title,

\begin{titlepage}
    \centering
    \vspace*{3cm}
    
    {\Huge CS3MI3 Assignment 3}\\[2cm]
    
    {\large Luca Mawyin}\\[0.5cm]
    
    {\large Dr.\ JP Rooney}\\[0.5cm]

    {\large \today}\\[2cm]

    {\Large McMaster University}

    \vfill
    \begin{flushleft}

        {\large \textbf{Student Number: }400531739}\\[0.5cm]
        {\large \textbf{MacID:} mawyinl}
    \end{flushleft}

\end{titlepage}

\renewcommand{\thesection}{Question \arabic{section}}
\renewcommand{\thesubsection}{\thesection.\alph{subsection}}

\titleformat{\subsection}
    {\normalfont\Large\bfseries}
    {\thesubsection}
    {1em}
    {}

\stepcounter{section}

\subsection{}

Firstly, we are given the definitions:
\begin{align*}
    c_0 &= \lambda{s}.~\lambda{z}.~z\\
    iszro &= \lambda{m}.~m~(\lambda{x}.~fls)~tru\\
    fls &= \lambda{t}.~\lambda{f}.~f\\
    tru &= \lambda{t}.~\lambda{f}.~t\\
\end{align*}
In order to evaluate 
\begin{align*}
    r &= iszro~c_0
\end{align*}
We will first apply $iszro$ to $c_0$ and evaluate:
\begin{align*}
    iszro~c_0 &\to (\lambda{m}.~m~(\lambda{x}.~fls)~tru)~c_0\\
    &\to c_0~(\lambda{x}.~fls)~tru
\end{align*}
Then we will substitute $c_0$ and evaluate:
\begin{align*}
    c_0~(\lambda{x}.~fls)~tru &\to (\lambda{s}.~\lambda{z}.~z)~(\lambda{x}.~fls)~tru\\
    &\to (\lambda{z}.~z)~tru\\
    &\to tru
\end{align*}
Therefore $r = iszro~c_0$ evaluates to $r=tru$.

\subsection{}

Firstly, we are given the definitions:
\begin{align*}
    or &= \lambda{b}.~\lambda{c}.~b~tru~c\\
    fls &= \lambda{t}.~\lambda{f}.~f\\
    tru &= \lambda{t}.~\lambda{f}.~t\\
\end{align*}
In order to evaluate 
\begin{align*}
    s &= or~tru~fls
\end{align*}
Firstly we will substitute $or$ and evaluate:
\begin{align*}
    or~tru~fls &\to (\lambda{b}.~\lambda{c}.~b~tru~c)~tru~fls\\
    &\to (\lambda{c}.~tru~tru~c)~fls\\
    &\to tru~tru~fls
\end{align*}
Then we will substitute $tru$ and evaluate:
\begin{align*}
    tru~tru~fls &\to (\lambda{t}.~\lambda{f}.~t)~tru~fls\\
    &\to (\lambda{f}.~tru)~fls\\
    &\to tru
\end{align*}
Therefore $s = or~tru~fls$ evalutates to $s = tru$.

\subsection{}

Firstly, we are given the definitions:
\begin{align*}
    and &= \lambda{b}.~\lambda{c}.~b~c~fls\\
    fls &= \lambda{t}.~\lambda{f}.~f\\
    tru &= \lambda{t}.~\lambda{f}.~t\\
\end{align*}
In order to evaluate 
\begin{align*}
    t &= and~tru~fls
\end{align*}
Firstly we will substitute $and$ and evaluate:
\begin{align*}
    and~tru~fls &\to (\lambda{b}.~\lambda{c}.~b~c~fls)~tru~fls\\
    &\to (\lambda{c}.~tru~c~fls)~fls\\
    &\to tru~fls~fls
\end{align*}
Then we will substitute $tru$ and evaluate:
\begin{align*}
    tru~fls~fls &\to (\lambda{t}.~\lambda{f}.~t)~fls~fls\\
    &\to (\lambda{f}.~fls)~fls\\
    &\to fls
\end{align*}
Therefore $t = and~tru~fls$ evaluates to $t = fls$.
\subsection{}

Using the results from the previous parts, we have definitions:
\begin{align*}
    r &= tru\\ 
    s &= tru\\
    t &= fls\\
    test &= \lambda{l}.~\lambda{m}.~\lambda{n}.~l~m~n
\end{align*}
In order to evaluate:
\begin{align*}
    test~r~s~t
\end{align*}
Firstly we will substitute $r$, $s$, and $t$:
\begin{align*}
    test~r~s~t &\to^* test~tru~s~t\\
    &\to^* test~tru~tru~t\\
    &\to^* test~tru~tru~fls
\end{align*}
Then we will expand $test$ and beta reduce:
\begin{align*}
    test~tru~tru~fls &\to (\lambda{l}.~\lambda{m}.~\lambda{n}.~l~m~n)~tru~tru~fls\\
    &\to (\lambda{m}.~\lambda{n}.~tru~m~n)~tru~fls\\
    &\to (\lambda{n}.~tru~tru~n)~fls\\
    &\to tru~tru~fls\\
    &\to (\lambda{f}.~tru)~fls\\
    &\to tru
\end{align*}
Therefore $test~r~s~t$ evaluates to $tru$.

\section{}

First we will show the rules for both languages. For arith we have:
\begin{center}
    $\frac{}{\text{if true then $t_2$ else $t_3\to{t_2}$}}$\\[1em]
    $\frac{}{\text{if false then $t_2$ else $t_3\to{t_3}$}}$\\[1em]
    $\frac{t_1\to{t'_1}}{\text{if true then $t_2$ else $t_3\to{t_2}$}\to\text{if $t'_1$ then $t_2$ else $t_3\to{t_2}$}}$
\end{center}
And in L-Bools we have:
\begin{align*}
    tru &= \lambda{m}.~\lambda{n}.~m\\
    fls &= \lambda{m}.~\lambda{n}.~n\\
    test &= \lambda{l}.~\lambda{m}.~\lambda{n}.~l~m~n
\end{align*}
Given the definitions, the evaluation relations do not agree. This is because arith will evaluate the condition before choosing the branch to return. This differs from L-Bools, as it will evaluate the branches first using a call-by-value evaluation that is grouped into a single step, before selecting one. We can demonstrate the difference using an example of an expression.\\[1em]
We will first represent it in arith:
\begin{align*}
    s &=\text{if false then false else }(\text{if true then true else false})\\
    &\to \text{if true then true else false}\\
    &\to \text{true}\\
\end{align*}
Then, in L-Bools we get:
\begin{align*}
    J(s) &= test~fls~fls~(test~tru~tru~fls)\\
    &\to test~fls~fls~tru\\
    &\to tru
\end{align*}
This shows that they do not agree. We can elaborate further by defining $s'$ for arith:
\begin{align*}
    s' &= \text{if true then true else false}\\
    s &\to s'
\end{align*}
and L-Bools:
\begin{align*}
    J(s') &= test~tru~tru~fls\\
    J(s) &\to tru
\end{align*}
As shown, arith evaluates the expression to the branch, while L-Bools evaluates the expression to $test$. Consequently, this shows that the evaluation strategies do not agree.

\end{document}